\section{Introduction}
A single coordinate sweep of the Thorp shuffle sends each card to a uniform position. This marginal statement leaves open the correlations among cards that have used the same switches. Those correlations are the object of this paper. We ask how much simultaneous row and column structure a permutation can retain, and how that restriction controls all the singular values of the sweep.

Let $n=2^d$ and identify the positions with $\{0,1\}^d$. A sweep visits the coordinates in order. At a visit it independently swaps, with probability $1/2$, the two cards on each edge in that coordinate direction. Write $T_n$ for its convolution operator on $\ell^2(S_n)$. One sweep is $d$ physical Thorp shuffles after removing the deterministic rotation of coordinates; Section~\ref{sec:prelim} gives the exact convention. The adjoint reverses the coordinate order. Thus $T_n^*T_n$ is a positive operator for a forward sweep followed by a reverse sweep, with fresh independent switches. Powers of $T_n$ instead describe repeated forward sweeps.

Thorp introduced the shuffle in 1973 while studying imperfect shuffling in Faro; his original description already makes the independent pair choices explicit~\cite[Section~3.1]{Thorp1973}. The distinction between one card and the whole deck is fundamental. One sweep uses $nd/2$ random bits, whereas a uniform permutation has entropy $\log_2(n!)$ bits. More generally, the support after $t$ physical shuffles has size at most $2^{tn/2}$, giving the lower bound $2d-O(1)$ for fixed-error mixing. Uniformity of every one-card marginal therefore leaves a substantial global problem.

The first polylogarithmic full-deck bound was Morris's $O(d^{44})$ theorem, announced in 2005 and published in 2008~\cite{Morris2008}. Montenegro and Tetali sharpened its evolving-set analysis to $O(d^{29})$~\cite[Section~6.4, Theorems~6.14--6.15]{MontenegroTetali2006}. Morris then introduced entropy contraction arguments giving $O((\log n)^4)$ shuffles for every even deck size~\cite{Morris2009}, and $O(d^3)$ for $n=2^d$~\cite{Morris2013}. These comparisons count physical shuffles at fixed total-variation error. The present paper obtains $\Theta(d)$ mixing for power-of-two decks from a stronger statement about the entire singular spectrum of one sweep. The upper bound uses a fixed number of complete sweeps.

Partial-card laws lead to a related but different question. Morris, Rogaway and Stegers obtained bounds for designated labeled cards and used them to analyze enciphering on small domains~\cite{mrs,MorrisRogawayStegers2018}. Czumaj and V\"ocking proved mixing of any fixed fraction less than one of the labeled cards after $O((\log n)^2)$ ordinary Thorp shuffles~\cite{CzumajVocking2014}. Their passage to a full random permutation uses additional indistinguishable cards which are discarded afterward. Czumaj also constructed switching networks that mix a full permutation with $O(n\log n)$ switches and $O((\log n)^2)$ depth~\cite{Czumaj2015}. These results explain the importance of shared-switch correlations beyond individual marginals; their state spaces or networks differ from the ordinary full-deck chain studied here.

Splitting the coordinate list in half identifies the positions with a rectangle. The first part of the sweep acts independently within its rows and the second independently within its columns. This is a recursion into smaller sweeps, but its two symmetry decompositions do not commute. We measure this obstruction by the following concrete event.

Take an $A$-by-$D$ rectangle with $AD=n$, and let $R=(S_D)^A$ and $C=(S_A)^D$ be its row and column permutation groups. A row permutation $r=(r_i)$ followed by a column permutation $c=(c_k)$ sends $(i,j)$ to $(c_{r_i(j)}(i),r_i(j))$. We call $(r,c)$ \emph{compatible} if this permutation can also be performed by a column permutation followed by a row permutation. Equivalently, for each original column $j$, the $A$ values $c_{r_i(j)}(i)$ are distinct. The opposite-order factorization, if it exists, is unique. Let $\mathcal I\subset R\times C$ denote this event and let $U$ denote independent uniform permutations in all the lines.

\begin{figure}[ht]
\centering
\begin{tikzpicture}[>=Stealth,scale=1.05]
\node (ij) at (0,1.8) {$(i,j)$};
\node (ik) at (4.2,1.8) {$(i,k)$};
\node (wj) at (0,0) {$(w,j)$};
\node (wk) at (4.2,0) {$(w,k)$};
\draw[->,thick,blue] (ij) -- node[above] {$r_i(j)=k$} (ik);
\draw[->,thick,red!70!black] (ik) -- node[right] {$c_k(i)=w$} (wk);
\draw[->,thick,red!70!black,dashed] (ij) -- node[left] {$\widetilde c_j(i)=w$} (wj);
\draw[->,thick,blue,dashed] (wj) -- node[below] {$\widetilde r_w(j)=k$} (wk);
\end{tikzpicture}

\caption{The two routings of one labeled position. The solid route uses a row permutation and then a column permutation. Compatibility means that the dashed routes of all positions can be assembled into column and row permutations; this is exactly the distinctness condition for the values $w$ in each original column $j$.}
\label{fig:compatibility}
\end{figure}

The first main result allows weights on every line. It retains the cost of marginal concentration instead of assuming the line laws remain uniform after compatibility is imposed.
\begin{theorem}[Weighted compatibility]\label{tra:duality:weighted}\label{tra:duality:weighted-statement}
There are absolute constants $L,C_0>0$ such that the following holds. Put $m=\sqrt n$, assume $m/2\le A,D\le2m$, and take $m$ sufficiently large that $\theta=1-L/\log m>0$. For all nonnegative functions $w_i:S_D\to[0,\infty)$ and $v_k:S_A\to[0,\infty)$,
\begin{equation}\label{tra:duality:eq5}
 \frac{n!}{|R||C|}\,
 \mathbb E_U\!\left[\mathbf1_{\mathcal I}
       \prod_{i=1}^A w_i(r_i)\prod_{k=1}^D v_k(c_k)\right]
 \le e^{C_0n^{.54}}
       \prod_{i=1}^A(\mathbb E w_i^{1/\theta})^\theta
       \prod_{k=1}^D(\mathbb E v_k^{1/\theta})^\theta .
\end{equation}
The expectations on the right are uniform on the indicated symmetric groups, and logarithms are natural.
\end{theorem}
The factor $n!/(|R||C|)$ is forced by regular trace normalization. Its leading logarithm is $n$. The entropy argument supplies a compatibility cost with leading logarithm $-n$, so those terms cancel. What remains is small enough to be absorbed by the multiplicity of a large irreducible representation.

To state the resulting spectral estimate, write $T_{n,\lambda}$ for the Fourier matrix on one copy of the irreducible representation of $S_n$ indexed by $\lambda\vdash n$, and let $D_\lambda$ be its dimension. Throughout, traces and Schatten norms are unnormalized. For $p>0$, the ordinary regular trace is
\[
 Z_n(p)=\Tr_{\mathrm{reg}}(T_n^*T_n)^p
       =\sum_{\lambda\vdash n}D_\lambda
                 \Tr(T_{n,\lambda}^*T_{n,\lambda})^p.
\]
The factor $D_\lambda$ records the number of copies in the regular representation.
\begin{theorem}[A bounded regular moment]\label{tra:duality:moment}
There is an absolute finite $p_*>0$ such that, for every dyadic $n$,
\begin{equation}\label{tra:duality:eq7}
 Z_n(p_*)\le1+\frac1{16}.
\end{equation}
For every integer $M\ge p_*$, $M$ independent forward sweeps have worst-start total-variation distance at most $1/8$ from uniform.
\end{theorem}
The moment records both the size and the number of nonconstant singular values. Its concrete consequences are as follows; Section~\ref{sec:spectral-conversion} proves the conversion with the regular multiplicities included.
\begin{corollary}[Singular ranks and full-deck mixing]\label{cor:intro-consequences}
For every nontrivial $\lambda\vdash n$ and $1\le r\le D_\lambda$,
\[
 s_r(T_{n,\lambda})\le(16D_\lambda r)^{-1/(2p_*)}.
\]
If $s_j(T_n)$ denotes the full regular singular list, including all multiplicities and zeros, then
\[
 s_j(T_n)\le j^{-1/(2p_*)}\qquad(1\le j\le n!).
\]
The constant singular value is the first value. The physical Thorp chain has worst-start $1/4$ mixing time $\Theta(d)$; more precisely its lower bound is $2d-O(1)$. For every fixed integer $M>p_*$, the worst-start total-variation distance after $Md$ physical shuffles tends to zero as $d\to\infty$.
\end{corollary}
The lower bound uses only support size. The upper bounds use the regular moment and Schatten H\"older for the matrix product $T_n^M$. They do not require the sweep to be normal. The sign block vanishes, and the dimensions of the remaining nonconstant blocks tend to infinity; the extra power when $M>p_*$ therefore makes the error vanish. The proof does not optimize the constant multiplying $d$ or establish a cutoff profile.

\paragraph{The first complete proof.}
Section~\ref{tra:duality} proves both theorems without using any of the later recursions. It starts with the exceptional representations, those whose Young diagram has a long first row. Such a representation first appears on ordered tuples of a corresponding number of cards. Subtracting all smaller-tuple kernels cancels every path configuration with an isolated card. The remaining configurations contain a forest of interactions. The probability of each forest edge, together with a deterministic bound on shared switches, gives the required sparse estimate.

For the other representations, the proof returns to the rectangle. Under an arbitrary law supported on compatible pairs, expose the row permutations and then the column permutations in a random order. Each light, unclumped entry costs nearly one nat because previously exposed columns forbid values. Repeated row images and large atoms of the local prediction are excluded from this calculation; their losses are charged to the corresponding marginal relative entropies. The resulting inequality still retains nearly the full sum of marginal deficits. Entropy duality then gives Theorem~\ref{tra:duality:weighted}.

The final step is analytic. A positive-matrix trace inequality reduces a singular moment to four alternating row and column factors. The identity coefficient of that product is exactly a weighted compatibility sum. Inverse Hausdorff--Young bounds its line weights by the smaller sweeps' moments. This first gives a subexponential bound on the regular trace. Each singular value in a large-dimensional irreducible occurs many times, so a small increase in the moment exponent makes that part of the trace small. The sparse estimate treats the remaining part. The exponent increases by $1+O(1/d)$ at a split into approximately $d/2$ coordinates; the product of these increases stays bounded.

\paragraph{The additional estimates.}
Sections~\ref{sec:prelim}--\ref{sec:spectral-conversion} contain the common conventions, the first complete proof, and the conversion to ranks and mixing. The subsequent arguments provide alternative proofs and several finer estimates. Some capped or weighted compatibility bounds already follow from Theorem~\ref{tra:duality:weighted}; their separate proofs offer different ways to expose the grid. The conditional correlation, exceptional-set, sparse representation and level estimates retain further structure, as described below.

The first group refines the entropy at the middle split. Section~\ref{tra:conditional} keeps conditional column deficits after the row array is known, together with projection ranks. Sections~\ref{tra:inverse-color}--\ref{tra:fractional-density} successively use inverse-color availability, survival-tilted predictions, and fractional integrability of coefficient densities. These formulations permit different changes of measure while retaining a quantitative charge for the concentration of the line laws. Section~\ref{tra:asymmetric} permits a very unequal rectangle; fixed chunks of coordinates provide its sparse estimate.

A second group keeps track of representation structure. Section~\ref{tra:projection-deletion} estimates the loss of diagram level under matching deletion and restriction to two halves. Section~\ref{tra:first-interactions} organizes the sparse cancellation by first interactions. Section~\ref{tra:harmonic-overlaps} lifts harmonic functions from smaller tuples and bounds both operator and Hilbert--Schmidt overlaps. Section~\ref{tra:regular-ranks} carries the full regular singular list through its recursion, whereas Section~\ref{tra:collision-series} retains positive level sums weighted by dimensions of the tail tableaux. The collision generating function in the latter argument preserves positivity before the missing regular multiplicity is restored.

The final three sections develop further ways to cross the middle grid. In Section~\ref{sec:convex-grid}, a convex weight in the diagram level and logarithmic dimension compensates for levels lost under restriction. In Section~\ref{sec:band-grid}, a fourth-moment estimate for singular bands transfers a profile without summing those bands prematurely. Finally, Section~\ref{sec:smoothed-grid} constructs a positive operator series that majorizes arbitrary matrix-coefficient squares by densities whose low-level restrictions can be controlled pointwise. The bound holds for each line permutation, so it remains valid after conditioning on a complete compatible grid while the auxiliary Bernoulli subsets are still sampled independently. The proof uses a tableau restriction estimate and finite-lattice association, and then telescopes the costs of smoothing along the exposure.

\paragraph{Methodological context.}
The rectangular geometry has a close precedent in H\aa stad's square lattice shuffle, which alternates independent uniform permutations within all rows and within all columns~\cite{Hastad2006,Hastad2016}. The corrected theorem gives full-permutation mixing in a constant number of such layers. In the Thorp recursion each line is acted on by a smaller sweep. Thus the middle-grid estimate has to accept nonuniform line densities and their Fourier blocks. This is the purpose of the marginal entropy terms in Theorem~\ref{tra:duality:weighted}.

Random-order exposure belongs to the entropy method used by Radhakrishnan in his proof of Br\'egman's theorem and by Linial and Luria for high-dimensional permutation counting~\cite{Radhakrishnan1997,LinialLuria2014}. Our exposure estimates keep relative entropies of nonuniform marginal laws; we prove those estimates here. The passage from those deficits to weighted product inequalities is the finite-space entropy--Brascamp--Lieb duality of Carlen and Cordero-Erausquin~\cite[Theorem~2.1]{CarlenCorderoErausquin2009}. The operator transfer uses the Araki--Lieb--Thirring inequality~\cite{araki1990,audenaert2007} and the classical inverse Hausdorff--Young inequality, in the compact-group normalization stated in~\cite[Lemma~5.1(2)]{garcia-cuerva-parcet2004}. We specify their normalization at use. The representation arguments use branching, Pieri, the hook-length formula and the content formula in the conventions of Sagan, Stanley and Vershik--Okounkov~\cite{Sagan2001,Stanley1999,VershikOkounkov2005}.

The sparse estimates also use an occupied-set product inequality preserved by fair swaps. We prove the needed two-site calculation directly. It is consistent with the stronger preservation of negative dependence under partial symmetrization proved by Borcea, Br\"and\'en and Liggett~\cite[Theorems~4.9 and~4.20]{BorceaBrandenLiggett2009}. The pointwise smoothing argument uses the finite distributive-lattice association theorem of Fortuin, Kasteleyn and Ginibre~\cite[Proposition~1]{FortuinKasteleynGinibre1971}; the required tableau monotonicities and the subsequent operator construction are established here.

The companion on routing densities~\cite[Lemma~3.2]{OpenAIRouting2026} proves the deterministic contact lemma, with contacts counted as shared switches; we also give an entropy proof of its local form when it enters the sparse estimate. Counting both participating endpoints doubles that bound. All sparse probability estimates and entropy-to-spectrum closures needed here are proved in this paper. The short companion \emph{Optimal-order mixing of the Thorp shuffle}~\cite[Theorem~1.1]{OpenAIOptimal2026} uses conditional marginals and an eight-block representation lift to prove optimal-order mixing of the full deck. Its mixing theorem is not an input to the compatibility arguments here.
